\scanpage{4}
\mspage{2}

Considérons un élément de

\[
 C_{\mathrm{st}}(I)\simeq C_{\mathrm{aff\,st}}(I_0)
\]

$\xi=\mathrm{cl}(X,\Sigma=(\Delta_i)_{i\in I})$. Alors $I_0$ hérite d'une \uline{structure polygonale} \inserted{(dépendant de $\xi$)} (on suppose $\operatorname{card} I_0\geq 3$) De plus, on trouve un polygône $\widetilde I_0(\xi)$ qui est un revêtement à deux feuillets de $I_0(\xi)$. Si on

\diagram{\begin{tikzcd}[ampersand replacement=\&]
\widetilde I_0(\xi)\arrow[d,phantom,"\downarrow" description]\\ I_0(\xi)
\end{tikzcd}}

\deleted{explicite} \inserted{traduit} la structure $(X,\Sigma=(\Delta_i)_{i\in I})$, en mettant à part $\Delta=\Delta_{i_0}$, à l'aide d'un disque $D=\mathrm{Déc}(X,\Delta)$, \deleted{alors} \inserted{et} d'un \uline{système \inserted{stable} de cordes} \inserted{$\widetilde\Sigma$} dans $D=\widetilde X$ indexé par $I_0$, on trouve \deleted{$\widetilde I_0(\xi)\unsure{\simeq}$ ens des sommets sur $\partial D$.}

\[
\begin{aligned}
\widetilde I_0(\xi)&\simeq\text{ens des sommets sur }\partial D=\partial\widetilde X\\
&\qquad\text{(rev. à deux feuillets de }\Delta\text{)}
\end{aligned}
\]

Soit maintenant $S$ un ens. polygonal de \unsure{sommets|sommet|même} d'ordre $2n$, donné une fois pour toutes. On s'intéresse aux structures

\[
\begin{aligned}
&\unsure{\mathfrak s|\text{de}}\qquad\text{syst. stables de cordes }S\text{-épinglés}\\
&\hspace{4em}\text{(mod }S\text{-isom)}
\end{aligned}
\]

soit $\widetilde C_{\mathrm{st}}(S)$ l'ens de ces structures. Posons

\[
\deleted{\unsure{\mathfrak S}}\quad I_0=S/a_S\qquad a_S\ \text{antipodisme de }S
\]

On a une application naturelle

\[
\widetilde C_{\mathrm{st}}(S)\longrightarrow C_{\mathrm{aff\,st}}(I_0)\deleted{\unsure{(I_0)}}
\]

\deleted{dont la fibre en $\xi$ est en corr. 1--1 avec l'ens des \unsure{L}}
% archive sheet number: 3
